\documentclass[12pt,a4paper]{article}

\usepackage[T1]{fontenc}
\usepackage[utf8]{inputenc}
\usepackage[brazil]{babel}
\usepackage[a4paper,margin=2.5cm]{geometry}
\usepackage{booktabs}
\usepackage{tabularx}
\usepackage{pdflscape}
\usepackage{ragged2e}
\usepackage{microtype}
\usepackage{xcolor}
\usepackage{tcolorbox}
\usepackage[hidelinks]{hyperref}

\newcolumntype{Y}{>{\RaggedRight\arraybackslash}X}
\setlength{\parindent}{1.25cm}
\setlength{\parskip}{0.4em}
\renewcommand{\arraystretch}{1.15}

\begin{document}

%======================= CAPA =======================
\begin{titlepage}
\centering
\vspace*{0.5cm}

{\large PONTIFÍCIA UNIVERSIDADE CATÓLICA DE MINAS GERAIS\par}
{\normalsize Instituto de Ciências Exatas e Informática --- ICEI\par}
{\normalsize Curso de Ciência de Dados e Inteligência Artificial\par}

\vspace{3.5cm}

{\Large\bfseries ESTRATÉGIA DE BUSCA BIBLIOGRÁFICA RASTREÁVEL\par}
\vspace{0.6cm}
{\large Construção, calibração e registro de uma expressão de busca\\
para o projeto de pesquisa da disciplina\par}

\vspace{3.5cm}

{\large\bfseries Pedro Dias Soares\par}
\vspace{0.3em}
{\large\bfseries Eduardo Velloso Landeiro\par}

\vspace{3.5cm}

\begin{flushright}
\begin{minipage}{0.52\textwidth}
\RaggedRight\normalsize
Exercício apresentado à disciplina Projeto em Ciência de Dados IV do curso de
Ciência de Dados e Inteligência Artificial da Pontifícia Universidade Católica
de Minas Gerais.

\vspace{0.8em}
Professor: Rodrigo Richard Gomes
\end{minipage}
\end{flushright}

\vfill
{\normalsize Belo Horizonte\par}
{\normalsize 2026\par}
\end{titlepage}

\tableofcontents
\newpage

%======================= 1 =======================
\section{Objeto deste exercício}

Este documento registra a construção de uma estratégia de busca bibliográfica para o
projeto de pesquisa da disciplina, seguindo o procedimento da oficina: decomposição da
pergunta em blocos conceituais, montagem de um conjunto-controle, iteração da expressão
com diagnóstico documentado a cada ajuste, tradução para mais de uma base e registro
rastreável de todas as consultas.

Todas as consultas relatadas foram \textbf{efetivamente executadas} em 17 de setembro de
2026. Nenhum número apresentado neste documento foi estimado. O registro completo, com a
expressão exata de cada consulta, está no arquivo \texttt{log-busca.csv}, versionado no
repositório do projeto, e o código que reproduz as consultas está em \texttt{buscas.py}.

\begin{tcolorbox}[colback=yellow!8,colframe=yellow!55!black,title=\textbf{Consultas ainda não executadas}]
As bases \textbf{Scopus}, \textbf{ACM Digital Library} e \textbf{IEEE Xplore} exigem
assinatura institucional pelo Portal CAPES (CAFe). Suas expressões estão traduzidas na
Seção~\ref{sec:traducao}, mas \textbf{ainda não foram executadas}, e as linhas
correspondentes do registro estão marcadas como \texttt{PENDENTE}. Nenhum número foi
atribuído a elas.
\end{tcolorbox}

%======================= 2 =======================
\section{Passo 1 --- Pergunta de pesquisa}

A pergunta é a do projeto em curso, registrada na Ficha Inicial:

\begin{quote}
\itshape
Quais fatores estão associados à defasagem de disponibilidade regional de modelos de IA
generativa comercial de provedores estadunidenses (Anthropic, OpenAI e Google), nas
plataformas AWS Bedrock, Google Vertex AI e Microsoft Foundry, entre abril de 2024 e
junho de 2026?
\end{quote}

O desfecho é a defasagem em dias até um modelo ficar disponível em uma região de nuvem,
tratada como tempo até evento com censura à direita.

%======================= 3 =======================
\section{Passo 2 --- Decomposição em blocos conceituais}

A pergunta foi decomposta em três conceitos. A apostila adverte que mais de quatro blocos
quase sempre zera a busca; três é o limite adotado.

\begin{table}[htbp]
\centering\small
\caption{Blocos conceituais}
\begin{tabularx}{\textwidth}{l l Y}
\toprule
\textbf{Bloco} & \textbf{Conceito} & \textbf{Papel na pergunta} \\
\midrule
B1 & Técnica & A tecnologia cujo acesso se mede: IA generativa e capacidade computacional \\
B2 & Desfecho & A grandeza medida: defasagem, desigualdade ou lacuna de acesso \\
B3 & Recorte & A dimensão comparativa: entre países, blocos ou regimes regulatórios \\
\bottomrule
\end{tabularx}
\end{table}

%======================= 4 =======================
\section{Passo 3 --- Conjunto-controle}

O conjunto-controle (Tabela~\ref{tab:ctrl}) é composto por cinco trabalhos cuja existência
foi verificada contra o
registro da agência mantenedora do identificador (Crossref, DataCite e API do arXiv) em
etapa anterior do projeto. Nenhum foi obtido por sugestão de modelo de linguagem.

\begin{table}[htbp]
\centering\footnotesize
\caption{Conjunto-controle --- versão final}
\label{tab:ctrl}
\begin{tabularx}{\textwidth}{l Y l l}
\toprule
\textbf{Referência} & \textbf{Título} & \textbf{Veículo} & \textbf{DOI} \\
\midrule
Lutz (2019) & Digital inequalities in the age of artificial intelligence and big data &
\textit{Human Behavior and Emerging Technologies} & \texttt{10.1002/hbe2.140} \\
Ahmed \& Wahed (2020) & The de-democratization of AI: deep learning and the compute divide
in AI research & arXiv (preprint, cs.CY) & \texttt{10.48550/arXiv.2010.15581} \\
Mannuru et al.\ (2023) & Artificial intelligence in developing countries: the impact of
generative AI & \textit{Information Development} & \texttt{10.1177/02666669231200628} \\
Wang et al.\ (2024) & The artificial intelligence divide: who is the most vulnerable? &
\textit{New Media \& Society} & \texttt{10.1177/14614448241232345} \\
Sastry et al.\ (2024) & Computing power and the governance of artificial intelligence &
arXiv (preprint, cs.CY) & \texttt{10.48550/arXiv.2402.08797} \\
\bottomrule
\end{tabularx}
\end{table}

\subsection{Revisão do conjunto-controle durante a execução}

O conjunto-controle inicial incluía Knez (2023), \textit{Technology diffusion and uneven
development}. O diagnóstico da Seção~\ref{sec:diag} mostrou que esse trabalho não é
recuperável por nenhuma expressão que contenha um bloco de IA, porque ele trata de difusão
de tecnologias industriais e de diferenças salariais entre países, sem mencionar
inteligência artificial. Exigir sua recuperação obrigava a incluir o termo genérico
\texttt{technology} no bloco B1, o que elevou o retorno a 186.169 registros.

Knez (2023) foi portanto \textbf{removido do conjunto-controle} --- ele não deveria
aparecer nesta busca --- e substituído por Lutz (2019). O trabalho permanece na
fundamentação do projeto como âncora da literatura de difusão, recuperado por uma linha de
busca distinta e declarada.

%======================= 5 =======================
\section{Passo 4 --- Sinônimos por bloco}

Os sinônimos não foram inventados. Foram extraídos dos \textbf{resumos efetivamente
recuperados} dos cinco trabalhos do conjunto-controle, conforme instrui a apostila.

\begin{table}[htbp]
\centering\footnotesize
\caption{Termo-âncora e variantes por bloco, com a origem de cada termo}
\begin{tabularx}{\textwidth}{l l Y}
\toprule
\textbf{Bloco} & \textbf{Termo-âncora} & \textbf{Variantes, siglas e origem} \\
\midrule
B1 & \texttt{artificial intelligence} &
\texttt{generative AI}, \texttt{generative artificial intelligence} (Mannuru);
\texttt{computing power}, \texttt{compute} (Sastry);
\texttt{large language model} (LLM), \texttt{foundation model};
\texttt{cloud computing} \\
\addlinespace
B2 & \texttt{digital divide} &
\texttt{digital inequality} / \texttt{inequalities} (Lutz);
\texttt{compute divide}, \texttt{de-democratization} (Ahmed \& Wahed);
\texttt{artificial intelligence divide} (Wang);
\texttt{unequal access} (Ahmed \& Wahed); \texttt{access gap};
\texttt{technology adoption}, \texttt{technology diffusion}, \texttt{adoption lag} \\
\addlinespace
B3 & \texttt{developing countries} &
\texttt{global south}; \texttt{cross-country}, \texttt{country differences} (Knez);
\texttt{digital sovereignty}; \texttt{AI governance}, \texttt{compute governance},
\texttt{governance of artificial intelligence} (Sastry) \\
\bottomrule
\end{tabularx}
\end{table}

\noindent\textbf{Vocabulário controlado.} Os vocabulários citados na literatura de origem
do método (MeSH, DeCS, EMTREE) pertencem à área da saúde. Os equivalentes em Computação
são o \textit{ACM Computing Classification System} e o \textit{IEEE Thesaurus}, adotados
como referência terminológica.

%======================= 6 =======================
\section{Passo 5 --- Versão 1 da expressão}

Aplicando a regra dos blocos --- \texttt{OR} dentro, \texttt{AND} entre, tudo entre
parênteses --- e traduzindo a pergunta literalmente:

\begin{quote}\footnotesize\ttfamily
("generative AI" OR "generative artificial intelligence" OR "large language model"\\
 OR "foundation model")\\
AND ("regional availability" OR "deployment delay" OR "adoption lag")\\
AND ("cloud region" OR "data residency" OR "digital sovereignty")
\end{quote}

\noindent\textbf{Resultado: zero registros.} A apostila classifica o retorno nulo como
sinal de alerta grave. Seguindo a orientação, a expressão foi verificada isolando um bloco
por vez: cada bloco isoladamente devolve resultados, o que descarta erro de sintaxe. O
diagnóstico é que os blocos são específicos demais, e que a interseção exata do objeto do
projeto é inédita na literatura indexada.

%======================= 7 =======================
\section{Passo 6 --- Iteração, diagnóstico e calibração}
\label{sec:diag}

A base de calibração foi a \textbf{OpenAlex}, escolhida por ser a única base indexadora de
acesso irrestrito disponível ao grupo. A calibração nas bases assinadas está pendente,
conforme declarado na abertura.

\subsection{Diagnóstico bloco a bloco}

Diante de três versões consecutivas sem recuperar nenhum artigo-controle, cada bloco foi
testado isoladamente contra cada artigo, em quinze consultas. Este é o diagnóstico que
reorientou todo o restante do trabalho.

\begin{table}[htbp]
\centering\small
\caption{Recuperação de cada artigo-controle por bloco isolado}
\begin{tabularx}{\textwidth}{Y c c c Y}
\toprule
\textbf{Artigo} & \textbf{B1} & \textbf{B2} & \textbf{B3} & \textbf{Por que falha} \\
\midrule
Ahmed \& Wahed (2020) & sim & sim & \textbf{não} & Unidade é a firma e a universidade, não o país \\
Knez (2023)           & \textbf{não} & sim & sim & Difusão de tecnologia em geral; não menciona IA \\
Wang et al.\ (2024)   & sim & sim & \textbf{não} & Survey com cidadãos holandeses; sem eixo entre países \\
Mannuru et al.\ (2023)& sim & \textbf{não} & sim & Trata de \textit{impacto}, sem termo de defasagem \\
Sastry et al.\ (2024) & sim & \textbf{não} & \textbf{não} & Governança de \textit{compute}; vocabulário próprio \\
\bottomrule
\end{tabularx}
\end{table}

\noindent Nenhum bloco está errado. O que a tabela mostra é que a estrutura de três blocos
em \texttt{AND} é estrita demais para um conjunto-controle que atravessa três literaturas
distintas. Duas consequências foram tiradas: a revisão do conjunto-controle descrita na
Seção~4.1, e o abandono da expressão única.

\subsection{Versões executadas}

\begin{table}[htbp]
\centering\footnotesize
\caption{Iteração da expressão na OpenAlex, 17 de setembro de 2026}
\begin{tabularx}{\textwidth}{l Y r c}
\toprule
\textbf{Versão} & \textbf{Ajuste aplicado e motivo} & \textbf{Result.} & \textbf{Controle} \\
\midrule
\multicolumn{4}{l}{\textit{Conjunto-controle inicial}}\\
v1 & Tradução literal da pergunta, três blocos estreitos & 0 & 0/5 \\
v2 & Sinônimos dos resumos acrescentados aos três blocos & 619 & 0/5 \\
v3 & B2 e B3 ampliados & 706 & 0/5 \\
\midrule
\multicolumn{4}{l}{\textit{Após o diagnóstico e a revisão do conjunto-controle}}\\
v4 & Dois blocos; \texttt{technology} genérico em B1 & 186.169 & 4/5 \\
v5 & Dois blocos; B1 restrito a IA & 31.285 & 4/5 \\
v6 & B1 \texttt{AND} B2 (família \textit{divide}) & 7.638 & 3/5 \\
v7 & v6 mais duas variantes de B2 & 7.642 & 3/5 \\
v8 & B1 \texttt{AND} (B2 \texttt{OR} B3) & 31.857 & \textbf{5/5} \\
v9 & B1 \texttt{AND} B2 \texttt{AND} B3 & 681 & 0/5 \\
v10 & v8 com filtro de ano ($\geq$ 2019) & 25.501 & \textbf{5/5} \\
v11 & v10 restrito a artigo e \textit{preprint} & 18.219 & \textbf{5/5} \\
v12 & Bloco de desfecho restrito ao título & 14.385 & \textbf{5/5} \\
v13 & Ambos os conceitos exigidos no título & 18.419 & \textbf{5/5} \\
\bottomrule
\end{tabularx}
\end{table}

\subsection{O impasse e sua resolução}

As versões v8 a v13 recuperam o conjunto-controle completo, mas nenhuma devolve volume
triável. A v9, que é triável, não recupera nada. Precisão e revocação estão em oposição
direta, e a causa é identificável: Mannuru e Sastry só são alcançados por termos de
território e governança, ao passo que Lutz, Ahmed e Wang só o são por termos da família
\textit{divide}. Essas duas famílias pertencem a literaturas diferentes.

Antes de aceitar esse resultado, verificou-se se o volume não seria artefato de sintaxe.
A base foi testada quanto ao tratamento de frase exata: \texttt{"artificial intelligence"}
devolve 327.044 registros no título contra 154 da frase invertida, e \texttt{"compute
divide"} devolve 11 contra 1 --- a base respeita aspas. O volume reflete o tamanho real da
literatura, não erro de construção.

\medskip
\noindent A resolução adotada é \textbf{declarar duas linhas de busca}, cada uma triável e
cada uma correspondendo a uma literatura, cuja união satisfaz o critério de aceitação.

\begin{table}[htbp]
\centering\footnotesize
\caption{Expressões finais --- OpenAlex, 17/09/2026, publicações a partir de 2019, artigos e \textit{preprints}}
\begin{tabularx}{\textwidth}{l Y r c}
\toprule
\textbf{Linha} & \textbf{Bloco combinado com B1} & \textbf{Result.} & \textbf{Controle} \\
\midrule
A & Desigualdade de acesso (família \textit{divide}) & 5.219 & 3/5 \\
B & Território e governança & 14.107 & 2/5 \\
\midrule
\multicolumn{2}{l}{\textbf{União de A e B}} & --- & \textbf{5/5} \\
\bottomrule
\end{tabularx}
\end{table}

\noindent\textbf{Deriva entre medições.} As duas linhas foram medidas duas vezes no mesmo
dia, com algumas horas de intervalo: a Linha A passou de 5.214 para 5.219 registros e a
Linha B de 14.082 para 14.107. A recuperação do conjunto-controle não mudou. A variação é
pequena, mas confirma em medição própria a advertência da apostila de que bases crescem e
de que um número sem data não significa nada. Os valores da tabela são os da segunda
medição.

\begin{quote}\footnotesize\ttfamily
LINHA A =\\
("artificial intelligence" OR "generative AI" OR "large language model"\\
 OR "foundation model" OR "computing power")\\
AND ("digital divide" OR "digital inequality" OR "digital inequalities"\\
 OR "compute divide" OR "AI divide" OR "artificial intelligence divide"\\
 OR "unequal access" OR "access gap")

\vspace{0.6em}
LINHA B =\\
("artificial intelligence" OR "generative AI" OR "large language model"\\
 OR "foundation model" OR "computing power")\\
AND ("developing countries" OR "global south" OR "digital sovereignty"\\
 OR "AI governance" OR "compute governance"\\
 OR "governance of artificial intelligence")
\end{quote}

%======================= 8 =======================
\section{Passo 7 --- Tradução para as demais bases}
\label{sec:traducao}

A estrutura conceitual se mantém; a escrita muda. Cada expressão abaixo está na sintaxe da
base a que se destina.

\subsection{arXiv --- executada}

Prefixos de campo próprios (\texttt{ti:}, \texttt{abs:}, \texttt{cat:}).

\begin{quote}\footnotesize\ttfamily
(abs:"artificial intelligence" OR abs:"generative AI" OR abs:"computing power")\\
AND (abs:"digital divide" OR abs:"compute divide" OR abs:"unequal access"\\
 OR abs:"digital sovereignty" OR abs:"developing countries")
\end{quote}

\noindent\textbf{100 registros} em 17/09/2026. A v1 estreita, traduzida para a mesma base,
devolveu \textbf{zero} --- confirmação independente do diagnóstico da Seção~6.

\subsection{SBC OpenLib --- executada}

Biblioteca digital da Sociedade Brasileira de Computação. A busca integrada não aceita
operadores booleanos, de modo que a tradução é por termos em português.

\begin{table}[htbp]
\centering\footnotesize
\begin{tabularx}{\textwidth}{Y r}
\toprule
\textbf{Termos (busca integrada)} & \textbf{Result.} \\
\midrule
soberania digital inteligência artificial & 6 \\
desigualdade digital inteligência artificial & 5 \\
computação em nuvem disponibilidade região & 1 \\
difusão tecnológica inteligência artificial & 1 \\
\bottomrule
\end{tabularx}
\end{table}

\noindent O retorno de um único registro para o recorte mais próximo do objeto do projeto
é, em si, um resultado: a interseção exata praticamente não existe na produção brasileira
indexada pela SBC.

\subsection{Scopus --- a executar}

\begin{quote}\footnotesize\ttfamily
TITLE-ABS-KEY(\\
 ("artificial intelligence" OR "generative AI" OR "large language model"\\
  OR "foundation model" OR "computing power")\\
 AND ("digital divide" OR "digital inequality" OR "compute divide"\\
  OR "unequal access" OR "access gap")\\
) AND PUBYEAR > 2018
\end{quote}

\noindent Agrupar sempre com parênteses: a Elsevier alterou a precedência dos operadores
booleanos, com implantação concluída no início de 2026, e expressões antigas sem
parênteses podem retornar resultado diferente.

\subsection{ACM Digital Library --- a executar}

\begin{quote}\footnotesize\ttfamily
Title:("artificial intelligence" OR "generative AI" OR "computing power")\\
AND Abstract:("digital divide" OR "compute divide" OR "unequal access")
\end{quote}

\noindent Na busca avançada, termos na mesma linha são combinados com \texttt{OR} por
padrão e linhas diferentes com \texttt{AND}. Escrever os operadores explicitamente e
conferir em \textit{View Query Syntax} antes de anotar o número.

\subsection{IEEE Xplore --- a executar}

\begin{quote}\footnotesize\ttfamily
("Abstract":"artificial intelligence" OR "Abstract":"computing power")\\
AND ("Abstract":"digital divide" OR "Abstract":"compute divide")
\end{quote}

\noindent \textit{Command Search}. Truncamento exige no mínimo três caracteres; limite de
25 termos por cláusula; por padrão busca apenas metadados.

%======================= 9 =======================
\section{Passo 8 --- Registro consolidado}

O registro completo está em \texttt{log-busca.csv}, versionado no repositório do projeto,
com 27 linhas cobrindo base, data, expressão exata, campos, filtros, número de resultados,
recuperação do conjunto-controle e observação. As três linhas de bases assinadas estão
marcadas \texttt{PENDENTE}, sem número atribuído.

%======================= 10 =======================
\section{Limitações declaradas}

\textbf{Bases assinadas não consultadas.} Scopus, Web of Science, ACM Digital Library e
IEEE Xplore exigem acesso institucional e não foram executadas. A calibração foi feita na
OpenAlex, cujo escopo é mais amplo e menos curado, o que explica em parte o volume elevado.
A repetição na Scopus deve reduzir os números e é a próxima etapa.

\textbf{SciELO.} A busca direta não respondeu a partir do ambiente utilizado.

\textbf{Google Scholar.} Não foi usado como base de busca. A consulta é cortada em cerca de
256 caracteres sem aviso, o teto de exibição é de aproximadamente mil resultados e o escopo
não é declarado --- não é reprodutível, e portanto não serve a um registro rastreável.

\textbf{Uso de IA generativa.} Modelo de linguagem foi usado para sugerir variantes de
termo, uso previsto na apostila. Nenhuma referência foi obtida por sugestão de modelo: os
cinco trabalhos do conjunto-controle têm DOI verificado contra Crossref, DataCite e a API
do arXiv.

%======================= 11 =======================
\section{Link do projeto \texorpdfstring{\LaTeX}{LaTeX}}

\begin{tcolorbox}[colback=red!5,colframe=red!60!black,title=\textbf{PREENCHER ANTES DE ENTREGAR}]
\centering\large\ttfamily
[INSERIR AQUI O LINK DO OVERLEAF]
\end{tcolorbox}

\noindent O projeto deve estar configurado como \textit{Anyone with this link can view}.

\end{document}
